\section{RQ1.4 Multi-Microservice Kubernetes Inference Chain}
\label{sec:rq14}

\subsection{Research Question}

\begin{quote}
\textit{How does increasing the number of microservices affect end-to-end inference
cost when ResNet-18 is decomposed into synchronously chained coarse architectural stages
under in-cluster Kubernetes execution?}
\end{quote}

RQ1.4 moves from the local split-selection stages of RQ1.1--RQ1.3 to the final
thesis-facing local Kubernetes benchmark for the coarse-stage ResNet-18 chain family.
Rather than selecting a boundary through carry-forward, this stage benchmarks a
predefined set of five in-cluster conditions from \texttt{monolithic\_k8s\_1svc} to
\texttt{chain\_5svc} to measure how end-to-end latency changes as synchronous service
boundaries are added under controlled local Kubernetes execution.

The literature framing for multi-boundary chain cost and Kubernetes-level
deployment overhead is discussed in
\cref{sec:rw-split-computing,sec:rw-cloud-native}.

\subsection{Experimental Setup}

The thesis-facing RQ1.4 benchmark used the frozen controlled local Kubernetes run
exported on 2026-04-21. The environment snapshot records an Ubuntu/Linux runtime
(\texttt{Linux-6.17.0-22-generic-x86\_64-with-glibc2.41}), Python~3.10.20, and
PyTorch~2.11.0+cu130. Deployment metadata shows that the benchmark ran in namespace
\texttt{rq14} on a single-node \texttt{kind} cluster
(\texttt{kind-thesis-rq14}); for every condition, all service pods were scheduled on the
same node, \texttt{thesis-rq14-control-plane}. The evaluated model was pretrained
ResNet-18 with CPU-only FP32 execution and input shape
$1 \times 3 \times 224 \times 224$.

\begin{table}[ht]
\centering
\caption{RQ1.4 benchmark configuration}
\label{tab:rq14_config}
\begin{tabular}{p{0.34\linewidth}p{0.58\linewidth}}
\toprule
\textbf{Setting} & \textbf{Value} \\
\midrule
Runtime & Local Ubuntu/Linux Kubernetes benchmark \\
Cluster & Single-node \texttt{kind} cluster (\texttt{kind-thesis-rq14}) \\
Namespace & \texttt{rq14} \\
Conditions & Five predefined configurations from \texttt{monolithic\_k8s\_1svc} to \texttt{chain\_5svc} \\
Execution mode & Local in-cluster rolling orchestration over the predefined condition set \\
Model & ResNet-18 (pretrained) \\
Device & CPU \\
Precision & FP32 \\
Input shape & $1 \times 3 \times 224 \times 224$ \\
Rounds & 5 \\
Warmup iterations & 50 \\
Measured iterations & 200 \\
Cooldown between conditions & 5 seconds \\
Random seed & 42 \\
Communication & gRPC (in-cluster Kubernetes) \\
gRPC port & 50051 \\
Maximum message size & 16~MB \\
Service pod CPU request / limit & 1 core / 1 core \\
Service pod memory request / limit & 512~MiB / 512~MiB \\
Client CPU request / limit & 1 core / 1 core \\
Client memory request / limit & 1~GiB / 1~GiB \\
Image pull policy & \texttt{IfNotPresent} \\
Carry-forward & Not applicable; RQ1.4 benchmarks a predefined configuration set \\
Parity validation & Local segment validation and live gRPC chain validation; \texttt{num\_inputs} $= 5$, \texttt{atol} $= 1 \times 10^{-5}$ \\
Warmup calibration & Trailing window $= 10$, CV threshold $= 0.02$ \\
\bottomrule
\end{tabular}
\end{table}

CPU stabilisation controls were requested for the benchmark processes and applied where
the runtime allowed. Threading controls for PyTorch, OpenMP, MKL, and OpenBLAS were all
requested at one thread and observed at one thread. CPU affinity to one physical core
with SMT excluded was requested and applied; the observed affinity set was core~0.
Process priority \texttt{nice~-5} was requested and applied. Governor and turbo controls
were also requested, but writes to \texttt{/sys} failed because the containerised
environment exposed those paths as read-only. The detected governor was already
\texttt{performance}, and turbo was already disabled. ASLR remained at the detected
value~2 and was not modified.

Methodological validity was established before timing. For all five conditions, local
segment validation against the monolithic model and live gRPC chain validation against
the deployed services both passed with \texttt{max\_abs\_diff} $= 0.0$,
\texttt{atol} $= 1 \times 10^{-5}$, and \texttt{num\_inputs} $= 5$.

\subsection{Conditions}

Five conditions were evaluated in a rolling orchestration sweep:

\begin{itemize}
    \item \textbf{\texttt{monolithic\_k8s\_1svc}}: Full ResNet-18 inference in a single Kubernetes service. This is the Kubernetes-native monolithic baseline.
    \item \textbf{\texttt{chain\_2svc}}: Two services with a split after \texttt{layer2}. Cumulative transferred activation per request: 980~KB.
    \item \textbf{\texttt{chain\_3svc}}: Three services with splits after \texttt{layer1} and \texttt{layer3}. Cumulative transferred activation per request: 1568~KB.
    \item \textbf{\texttt{chain\_4svc}}: Four services with splits after \texttt{layer1}, \texttt{layer2}, and \texttt{layer3}. Cumulative transferred activation per request: 1960~KB.
    \item \textbf{\texttt{chain\_5svc}}: Five services with splits after \texttt{layer1}, \texttt{layer2}, \texttt{layer3}, and \texttt{layer4}. Cumulative transferred activation per request: 2058~KB.
\end{itemize}

Carry-forward is not applicable in RQ1.4. The frozen \texttt{carry\_forward.json}
states that carry-forward is defined only for the two-service split-selection experiments
(RQ1.1 and RQ1.2), whereas RQ1.4 benchmarks this predefined Kubernetes chain set.

\subsection{Results}

\begin{table}[ht]
\centering
\caption{RQ1.4 summary of mean latency and overhead vs.\ Kubernetes monolithic baseline}
\label{tab:rq14_results}
\begin{tabular}{lccccc}
\toprule
\textbf{Condition} & \textbf{Mean (ms)} & \textbf{Median (ms)} & \textbf{Std (ms)} & \textbf{p95 (ms)} & \textbf{Overhead vs.\ baseline} \\
\midrule
\texttt{monolithic\_k8s\_1svc} & 81.447 & 79.648 & 4.061 & 89.444 & --- \\
\texttt{chain\_2svc}           & 84.189 & 82.880 & 3.294 & 90.538 & +2.742 / +3.4\% \\
\texttt{chain\_3svc}           & 87.608 & 85.487 & 4.274 & 96.083 & +6.161 / +7.6\% \\
\texttt{chain\_4svc}           & 92.312 & 90.594 & 4.942 & 101.223 & +10.865 / +13.3\% \\
\texttt{chain\_5svc}           & 92.624 & 90.312 & 4.318 & 99.398 & +11.177 / +13.7\% \\
\bottomrule
\end{tabular}
\end{table}

\begin{figure}[htbp]
\centering
\resizebox{0.92\textwidth}{!}{%
\begin{tikzpicture}[
    font=\small,
    axis/.style={
        thick,
        draw=black!70,
        -{Latex[length=2mm]}
    },
    grid/.style={
        draw=black!12,
        thin
    },
    point/.style={
        circle,
        draw=blue!65!black,
        fill=blue!20,
        thick,
        minimum size=5.5pt,
        inner sep=0pt
    },
    pointlabel/.style={
        font=\scriptsize\bfseries,
        text=blue!65!black,
        align=center
    },
    xlabel/.style={
        font=\scriptsize,
        text=black!75,
        align=center
    },
    ylabel/.style={
        font=\small,
        text=black!75
    },
    transferlabel/.style={
        font=\scriptsize,
        text=black!60,
        align=center
    },
    note/.style={
        font=\scriptsize,
        text=black!65,
        align=center
    },
    errbar/.style={draw=black!80, semithick, line cap=round}
]

% y-axis: 76--98 ms at 0.28 plot-units per ms, widened from the bar-top range
% so the +/- 1 SD error bars fit while keeping the original figure height.
% y(lat) = (lat - 76) * 0.28.
\foreach \lat/\y in {76/0,78/0.56,80/1.12,82/1.68,84/2.24,86/2.80,88/3.36,90/3.92,92/4.48,94/5.04,96/5.60,98/6.16} {
    \draw[grid] (0,\y) -- (10,\y);
    \node[anchor=east, font=\scriptsize, text=black!65] at (-0.18,\y) {\lat};
}

\draw[axis] (0,0) -- (10.45,0);
\draw[axis] (0,0) -- (0,6.46);

\node[ylabel, rotate=90] at (-1.05,3.08) {Mean end-to-end latency (ms)};
\node[ylabel] at (5,-1.35) {Number of Kubernetes inference services};

\coordinate (p1) at (1,1.418);
\coordinate (p2) at (3,2.068);
\coordinate (p3) at (5,3.284);
\coordinate (p4) at (7,4.080);
\coordinate (p5) at (9,4.630);

% Error bars: mean +/- 1 SD of per-iteration latency (Std column of the
% Stage~K results table). Half-cap width 0.18.
\draw[errbar] (1,0.294) -- (1,2.543);
\draw[errbar] (0.82,0.294) -- (1.18,0.294);
\draw[errbar] (0.82,2.543) -- (1.18,2.543);
\draw[errbar] (3,1.081) -- (3,3.055);
\draw[errbar] (2.82,1.081) -- (3.18,1.081);
\draw[errbar] (2.82,3.055) -- (3.18,3.055);
\draw[errbar] (5,2.099) -- (5,4.469);
\draw[errbar] (4.82,2.099) -- (5.18,2.099);
\draw[errbar] (4.82,4.469) -- (5.18,4.469);
\draw[errbar] (7,2.980) -- (7,5.180);
\draw[errbar] (6.82,2.980) -- (7.18,2.980);
\draw[errbar] (6.82,5.180) -- (7.18,5.180);
\draw[errbar] (9,3.520) -- (9,5.740);
\draw[errbar] (8.82,3.520) -- (9.18,3.520);
\draw[errbar] (8.82,5.740) -- (9.18,5.740);

\draw[very thick, draw=blue!65!black]
    (p1) -- (p2) -- (p3) -- (p4) -- (p5);

\node[point] at (p1) {};
\node[point] at (p2) {};
\node[point] at (p3) {};
\node[point] at (p4) {};
\node[point] at (p5) {};

\node[pointlabel, anchor=south] at (1,2.66) {81.07 ms};
\node[pointlabel, anchor=south] at (3,3.18) {83.39 ms};
\node[pointlabel, anchor=south] at (5,4.59) {87.73 ms};
\node[pointlabel, anchor=south] at (7,5.30) {90.57 ms};
\node[pointlabel, anchor=south] at (9,5.86) {92.54 ms};

\node[xlabel] at (1,-0.45) {\texttt{1svc}\\monolithic};
\node[xlabel] at (3,-0.45) {\texttt{2svc}};
\node[xlabel] at (5,-0.45) {\texttt{3svc}};
\node[xlabel] at (7,-0.45) {\texttt{4svc}};
\node[xlabel] at (9,-0.45) {\texttt{5svc}};

\node[transferlabel] at (1,-0.95) {588 KiB};
\node[transferlabel] at (3,-0.95) {980 KiB};
\node[transferlabel] at (5,-0.95) {1568 KiB};
\node[transferlabel] at (7,-0.95) {1960 KiB};
\node[transferlabel] at (9,-0.95) {2058 KiB};

\node[transferlabel, anchor=east] at (-0.25,-0.95) {recorded\\tensor bytes};

\node[note, text width=10cm] at (5,-1.95)
    {Error bars show $\pm 1$\,SD of per-iteration latency. The ordering is monotonic and the trend approximately linear (slope $\approx 3.0$\,ms/service, $R^{2}\approx0.98$); adjacent gaps compare predefined configurations.};

\end{tikzpicture}

}
\caption{Mean end-to-end latency for the RQ1.4 local Kubernetes chain family. Latency increases as additional synchronous service boundaries are introduced, but the increase is not strictly linear. The smallest marginal increase occurs between the four-service and five-service chains, where the added late-stage boundary contributes comparatively little additional activation-transfer burden.}
\label{fig:rq14_kubernetes_chain_latency}
\end{figure}

\Cref{fig:rq14_kubernetes_chain_latency} visualises the same latency trend reported in
\cref{tab:rq14_results}, showing that most of the increase occurs before the
five-service configuration.

\begin{table}[ht]
\centering
\caption{RQ1.4 overhead and inferred non-compute/platform cost breakdown}
\label{tab:rq14_overhead}
\begin{tabular}{lcccc}
\toprule
\textbf{Condition} & \textbf{Overhead (ms)} & \textbf{Overhead (\%)} & \textbf{Cumul.\ Activation (KB)} & \textbf{Inferred Non-Compute / Platform (ms)} \\
\midrule
\texttt{chain\_2svc} & 2.742  & 3.4  & 980  & 6.285 \\
\texttt{chain\_3svc} & 6.161  & 7.6  & 1568 & 9.646 \\
\texttt{chain\_4svc} & 10.865 & 13.3 & 1960 & 12.231 \\
\texttt{chain\_5svc} & 11.177 & 13.7 & 2058 & 13.893 \\
\bottomrule
\end{tabular}
\end{table}

\begin{table}[ht]
\centering
\caption{RQ1.4 effect sizes relative to monolithic baseline}
\label{tab:rq14_effects}
\begin{tabular}{lccc}
\toprule
\textbf{Condition} & \textbf{Cohen's $d$} & \textbf{Interpretation} & \textbf{Mann-Whitney $p$} \\
\midrule
\texttt{chain\_2svc} & 0.742 & medium & $2.38 \times 10^{-81}$ \\
\texttt{chain\_3svc} & 1.478 & large  & $6.82 \times 10^{-150}$ \\
\texttt{chain\_4svc} & 2.402 & large  & $3.03 \times 10^{-282}$ \\
\texttt{chain\_5svc} & 2.667 & large  & $3.50 \times 10^{-291}$ \\
\bottomrule
\end{tabular}
\end{table}

\noindent
\textit{Methodological note.} Effect sizes are computed from pooled per-iteration data
($N = 1{,}000$ iterations per condition). With large $N$, Mann-Whitney $p$-values are
near-zero for any non-trivial difference and should not be over-interpreted as evidence of
strong practical significance. Cohen's $d$ provides a more informative measure of effect
magnitude. Cross-round consistency and parity validation are the primary indicators of
result stability and methodological validity.

\begin{table}[ht]
\centering
\caption{RQ1.4 cross-round consistency}
\label{tab:rq14_rounds}
\begin{tabular}{lcccc}
\toprule
\textbf{Condition} & \textbf{Grand Mean (ms)} & \textbf{Round Std (ms)} & \textbf{Round Range (ms)} & \textbf{Round CV} \\
\midrule
\texttt{monolithic\_k8s\_1svc} & 81.447 & 1.3193 & 3.0408 & 0.0162 \\
\texttt{chain\_2svc}           & 84.189 & 1.1876 & 3.0546 & 0.0141 \\
\texttt{chain\_3svc}           & 87.608 & 1.4250 & 3.7465 & 0.0163 \\
\texttt{chain\_4svc}           & 92.312 & 1.6107 & 3.7851 & 0.0174 \\
\texttt{chain\_5svc}           & 92.624 & 0.6942 & 1.7056 & 0.0075 \\
\bottomrule
\end{tabular}
\end{table}

The frozen controlled local Kubernetes run was stable across rounds. Round CVs remained low for all
five conditions: 0.0162 for \texttt{monolithic\_k8s\_1svc}, 0.0141 for
\texttt{chain\_2svc}, 0.0163 for \texttt{chain\_3svc}, 0.0174 for
\texttt{chain\_4svc}, and 0.0075 for \texttt{chain\_5svc}. Together with the parity
checks above, this supports treating the reported condition means as representative of
the local in-cluster benchmark.

The interpretation of these results and the answer to RQ1.4 are presented in
\cref{sec:analysis-rq14}.